\documentclass[11pt,a4paper]{article}
\usepackage{fontspec,amsmath,amssymb,amsthm,microtype}
\setmainfont{lmroman10-regular.otf}[BoldFont=lmroman10-bold.otf,ItalicFont=lmroman10-italic.otf,BoldItalicFont=lmroman10-bolditalic.otf]
\setmonofont{lmmono10-regular.otf}[Scale=MatchLowercase]
\usepackage[margin=25mm]{geometry}
\usepackage[unicode,hidelinks]{hyperref}
\setlength{\emergencystretch}{3em}
\setlength{\parskip}{0.4em}
\setlength{\parindent}{0pt}
\allowdisplaybreaks[2]
\raggedbottom
\title{\texorpdfstring{Sharp transformed Gaussian bounds for the repaired random Hessian correction}{Sharp transformed Gaussian bounds for the repaired random Hessian correction}}
\author{LangevinLean mathematical exposition}
\date{}
\begin{document}\maketitle
\noindent Source module: \texttt{Kinetic\allowbreak{}Langevin/\allowbreak{}Free\allowbreak{}OUCorrection\allowbreak{}Bound.\allowbreak{}lean}. The accompanying HTML pairs this mathematical exposition with the exact Lean statements and records the scope of the formalization.


The actual \ensuremath{\varphi}\ensuremath{{}_{2}} and \ensuremath{\varphi}\ensuremath{{}_{3}} stochastic Hessian terms are estimated together. Their exact \ensuremath{\psi}\ensuremath{{}_{2}} cancellation gives the sharp quadratic time coefficient. Gaussian truncation then supplies the source-paper random-Hessian contribution without discarding fresh-noise terms or assuming independence from the Hessian.


\section{\texorpdfstring{Prove the sharp first and second primitive estimates}{Prove the sharp first and second primitive estimates}}


If \(\|U_r\|\le C\) on \([0,t]\), then \(\|\int_0^sU_r\,dr\|\le Cs\). Integrating this genuine bound yields \[\left\|\int_0^t\int_0^sU_r\,dr\,ds\right\|\le\frac{Ct^2}{2}.\] The same estimate holds with the outer factor \(e^{-\gamma(t-s)}\) for \(\gamma\ge0\), since that scalar lies between zero and one. Lean first applies the constant integral-norm bound on each actual inner interval, then bounds the outer integral by the integrable scalar function \(Cs\). The exact scalar primitive gives the factor \(1/2\). Every interval endpoint and nonnegative damping condition is checked explicitly.



{\raggedright\noindent Lean declarations: \texttt{norm\_\allowbreak{}primitive\_\allowbreak{}le\_\allowbreak{}linear}, \texttt{norm\_\allowbreak{}double\_\allowbreak{}primitive\_\allowbreak{}le\_\allowbreak{}quadratic}, \texttt{norm\_\allowbreak{}damped\_\allowbreak{}primitive\_\allowbreak{}le\_\allowbreak{}quadratic}.\par}

\section{\texorpdfstring{Use the full stochastic cancellation in the actual phase metric}{Use the full stochastic cancellation in the actual phase metric}}


Let \(A\) be the actual continuous linear random-Hessian operator and suppose \(\|AR_0(r)\|_{L^2}\le C\). The complete correction identities give \(AC_1\) as a damped second primitive of \(AR_0\), so \(\|AC_1(t)\|\le Ct^2/2\). The exact identity \(C_1+\gamma C_2=R_2\), including both stochastic responses, gives \(\|AC_1+\gamma AC_2\|=\|AR_2\|\le Ct^2/2\). The transformed phase norm combines exactly these two coordinates, hence \[\|T_{\gamma,0}(AC_1,AC_2)\|\le\frac{\sqrt2}{2}t^2C.\] Lean rewrites the two actual primitive identities, applies the two sharp integral bounds, commutes \(A\) with the exact correction balance, and sums the two squared component bounds. It finally checks the square of the displayed nonnegative bound. This controls the entire implemented correction.



{\raggedright\noindent Lean declarations: \texttt{free\allowbreak{}OUCorrection\_\allowbreak{}transform\allowbreak{}Norm\_\allowbreak{}bound}.\par}

\section{\texorpdfstring{Insert the actual Gaussian cutoff at each free-\allowbreak{}OU time}{Insert the actual Gaussian cutoff at each free-OU time}}


If every actual free-OU velocity marginal is standard Gaussian and \(\|AR_0(r)(\omega)\|\le\min\{B,M_2\|X(\omega)\|\}\|R_0(r)(\omega)\|\), the proved truncation theorem gives the same bound at every time: \[\|AR_0(r)\|_{L^2}\le M_2\sqrt\alpha\,\|X\|_{L^2}+2B e^{-(\alpha-p)/8},\qquad\alpha\ge5p.\] Substituting this common \(C\) in the sharp transformed estimate retains both terms with their exact constants. Lean constructs the nonnegative common bound and applies the actual Gaussian truncation theorem at each time. No independence between the Hessian factor, the position discrepancy, and the velocity is required.



{\raggedright\noindent Lean declarations: \texttt{free\allowbreak{}OUCorrection\_\allowbreak{}gaussian\_\allowbreak{}transform\allowbreak{}Norm\_\allowbreak{}bound}.\par}

\section{\texorpdfstring{Verify the cutoff premise for a concrete random operator}{Verify the cutoff premise for a concrete random operator}}


For an actual measurable random operator field \(H(\omega)\), suppose its true operator norm satisfies both \(\|H(\omega)\|\le B\) and \(\|H(\omega)\|\le M_2\|X(\omega)\|\) almost everywhere. The checked random-operator construction gives an actual L\ensuremath{{}^{2}} continuous linear map with representative \(H(\omega)V(\omega)\). Lean rewrites that actual representative, applies the operator norm inequality, combines the two operator bounds by their minimum, and invokes Gaussian truncation for a velocity with the actual standard Gaussian law. Thus the required cutoff estimate follows from the genuine uniform and Lipschitz operator bounds of a Hessian difference.



{\raggedright\noindent Lean declarations: \texttt{random\allowbreak{}Operator\allowbreak{}Lp\_\allowbreak{}gaussian\_\allowbreak{}truncation}.\par}

\end{document}
